\chapter{General Context of the Project}
\label{ch:context}

\vcyclephase{Project Context and Feasibility Analysis}

\section{Host Company / Facility Description}
\label{sec:company}

\subsection{General Company Overview}

ECI (Electrical Components International) was founded in 1953 as Burcliff Industries. It has become one of the world's leading suppliers of electrical distribution systems, control box assemblies, and critical technical components. The company currently employs 17,000 people across 30 manufacturing sites worldwide, serving over 450 clients. ECI designs, manufactures, and assembles electrical distribution systems, control boxes, and other critical technical components that enable the deployment of advanced technologies across various industrial sectors. Its slogan reflects this mission: \textit{"Powering Smart, Connected, and Electrified Solutions."}

\begin{figure}[H]
  \centering
  \textit{[Figure placeholder: Global locations of the ECI group]}
  \caption{Global locations of ECI.}
  \label{fig:eci-global}
\end{figure}

\subsection{Markets, Customers, and Suppliers}

ECI serves over 450 market-leading customers across a wide range of industries, positioning itself as a supplier capable of managing highly complex requirements in terms of products, demand, manufacturing, and delivery. Its main markets include:

\begin{itemize}
  \item \textbf{Home Appliance:} Global leader in wire harnesses and control box assemblies for major home appliance OEMs worldwide since 1953.
  \item \textbf{Automotive Advanced Safety:} Over 27 years of experience in manufacturing wire harnesses for ADAS systems, including cameras, sensors, and exterior mirrors.
  \item \textbf{Commercial Electric Vehicles:} Through its ProEVT brand, ECI leads the electrification of school buses, delivery trucks, and urban transport vehicles.
  \item \textbf{Aerospace:} Supply of wiring harnesses and kits for commercial, business, and military aircraft platforms.
  \item \textbf{Medical Devices:} Electrical distribution systems for diagnostic equipment, including MRI machines, X-ray imaging, and surgical applications.
  \item \textbf{Smart Industrial \& Access Equipment:} Wiring harnesses for forklifts, cranes, scissor lifts, and automated industrial machinery.
\end{itemize}

\begin{figure}[H]
  \centering
  \textit{[Figure placeholder: ECI suppliers]}
  \caption{ECI suppliers.}
  \label{fig:eci-suppliers}
\end{figure}

\begin{figure}[H]
  \centering
  \textit{[Figure placeholder: ECI customers]}
  \caption{ECI customers.}
  \label{fig:eci-customers}
\end{figure}

\subsection{ECI Subsidiaries}

\begin{figure}[H]
  \centering
  \textit{[Figure placeholder: ECI subsidiaries]}
  \caption{ECI subsidiaries.}
  \label{fig:eci-subsidiaries}
\end{figure}

\subsection{Vision and Values}

ECI's vision is to power the next generation of smart, connected, and electrified technologies. The company positions itself not merely as a component supplier, but as a long-term design and engineering partner, supporting its customers from concept development to large-scale production. Its four core values are:

\begin{itemize}
  \item \textbf{Operational Excellence:} A culture of \textit{zero defects}, Lean Manufacturing, and continuous improvement at all production levels.
  \item \textbf{Customer Commitment:} On-time delivery and first-class service represent ECI's absolute operational priority, supported by advanced demand analytics.
  \item \textbf{Ethical Excellence:} Integrity, collective responsibility, and commitment to promises at all levels of the organization.
  \item \textbf{Leadership in Innovation and Technology:} Continuous investment in electrification, automation, and digital manufacturing.
\end{itemize}

\subsection{ECI Morocco}

The Tangier plant, established in 2014 as the first investor in Tanger Automotive City (TAC), reached a major milestone in 2024 with the inauguration of its production site extension. This expansion increased annual production capacity to 40 million units and consolidated a workforce of 1,400 employees. This development underscores ECI's strategic commitment to Morocco and its active contribution to the industrial development of the Tangier-T\'{e}touan-Al Hoceima region.

\begin{figure}[H]
  \centering
  \textit{[Figure placeholder: ECI Morocco site]}
  \caption{ECI Morocco.}
  \label{fig:eci-morocco}
\end{figure}

\subsection{Company Profile Sheet}

\begin{table}[H]
  \centering
  \caption{Company profile sheet for ECI Tangier.}
  \label{tab:company-profile}
  \begin{tabular}{p{4cm} p{8cm}}
    \toprule
    \textbf{Attribute} & \textbf{Details} \\
    \midrule
    Company Name & Electrical Components International Morocco (ECI Morocco) \\
    Legal Form & Public Limited Company (PLC) \\
    Activity Type & Manufacturing and assembly of wire harnesses, electrical distribution systems and electronic components \\
    Year of Establishment & 2014 \\
    Capital & 166,099,000.00 MAD \\
    Workforce (Morocco) & Approximately 1,400 direct employees \\
    Production Capacity & 40 million units per year (following the 2024 expansion) \\
    Address & Zone Franche Tanger Automotive City (TAC), Tangier, Morocco \\
    Telephone & +212 531 063 615 \\
    Website & \texttt{www.ecintl.com} \\
    Parent Company & Electrical Components International, Inc. --- Headquarters: Southfield, Michigan, USA \\
    Certifications & ISO 9001:2015, IATF 16949, ISO 14001:2015 \\
    Markets Served & Automotive, Home Appliance, Commercial Electric Vehicles, Aerospace, Industrial Equipment \\
    Industrial Zone & Export Free Zone --- Tanger Med Industrial Platform \\
    \bottomrule
  \end{tabular}
\end{table}

\section{Production Zones}
\label{sec:production-zones}

The production flow at ECI Tangier is organized into four main zones:

\subsection{Warehouse and Feeding}

This is the starting point located on the left side of the layout. The warehouse stores raw materials that feed the various work areas. This is where the Pick-to-Light project intervenes to optimize component preparation before they are sent to production.

\subsection{Cutting and Preparation Zone}

Once materials leave the warehouse, they go through the Cutting Zone (APU Cutting). The operations performed include:

\begin{itemize}
  \item Cutting to length.
  \item Wire stripping and crimping.
  \item Connector insertion (Semi-IDC) and twisting.
\end{itemize}

\subsection{Assembly Zone}

This is the central and largest area of the plant. It consists of individual assembly stations where the final tasks are performed:

\begin{itemize}
  \item Taping (harness protection).
  \item Packaging and continuity or quality testing.
\end{itemize}

\subsection{Inspection and Shipping}

Products pass through a final control zone (no technical added value but crucial for quality) before reaching the shipping zone to be sent to customers.

\begin{figure}[H]
  \centering
  \textit{[Figure placeholder: General plant layout]}
  \caption{General plant layout of ECI Tangier.}
  \label{fig:plant-layout}
\end{figure}

\begin{figure}[H]
  \centering
  \textit{[Figure placeholder: Production flow diagram]}
  \caption{Production flow diagram.}
  \label{fig:production-flow}
\end{figure}

\subsection{Product Range}

\begin{figure}[H]
  \centering
  \textit{[Figure placeholder: Product range segmentation]}
  \caption{Product portfolio segmentation of ECI Morocco.}
  \label{fig:product-range}
\end{figure}

\subsection{Organizational Chart}

\begin{figure}[H]
  \centering
  \textit{[Figure placeholder: ECI Tangier organizational chart]}
  \caption{Departments of ECI Tangier.}
  \label{fig:org-chart}
\end{figure}

ECI Tangier is organized into complementary functional departments, each contributing to the overall operational efficiency of the plant. The general structure is presented in Table~\ref{tab:eci-departments}.

\begin{table}[H]
  \centering
  \caption{ECI departments.}
  \label{tab:eci-departments}
  \begin{tabular}{p{4cm} p{8cm}}
    \toprule
    \textbf{Department} & \textbf{Main Responsibilities} \\
    \midrule
    General Management & Overall management of the plant, strategy and performance supervision \\
    Production & Manufacture of cable harnesses, management of assembly lines, and achievement of performance objectives \\
    Quality Assurance & Product inspection, compliance with customer standards and defect prevention \\
    Process Engineering & Optimization of processes and digitalization initiatives --- Internship host department \\
    Supply Chain \& Warehouse & Raw material supply, inventory management and logistics \\
    Human Resources & Recruitment, training, and labor relations \\
    Finance \& Administration & Budgeting, reporting, and administrative operations \\
    NPI Product & New product introduction study \\
    \bottomrule
  \end{tabular}
\end{table}

\section{Project Context}
\label{sec:project-context}

\subsection{Pedagogical Context}

This project is part of the end-of-studies internship as a Master's student in Mechatronics at the Faculty of Sciences of Tetouan (FST). The main objective of this internship is to integrate into a professional environment and to develop the strengths acquired during the training by working in real situations.

\subsection{Educational General Context}

The Smart Storage Management System (SSMS) is a Pick-to-Light / Put-to-Light warehouse management solution designed for a partner industrial facility. The system uses WS2813 addressable RGB LEDs installed on storage bins to guide operators during picking and stocking operations, reducing human error and improving retrieval speed. The system is controlled by an ESP32-WROOM-32D microcontroller running MicroPython firmware, communicating via MQTT over a local LAN with a Windows desktop application built in Python/Tkinter.

\subsection{Process Engineering Department}

The Process Engineering Department at ECI Tangier is responsible for analyzing, designing, and continuously improving the plant's manufacturing and logistics processes. Operating at the intersection of industrial engineering and digital technologies, it constitutes the main driver of the company's Industry 4.0 strategy. It is within this department that this final year project takes place.

\section{Warehouse Overview}
\label{sec:warehouse}

ECI Tangier's raw materials warehouse manages the reception, storage, and supply of all components required for cable harness manufacturing: cables of various gauges, connectors, terminals, and protective sheaths.

\subsection{Physical and Logical Organization}

The warehouse is structured around 7 storage racks, subdivided into individual locations identified by a position code.

\subsection{Current Capacity}

25 active storage positions spread across 7 racks.

\subsection{State Management}

Each location is associated with a specific reference and has a binary status: FULL (busy) or EMPTY (available).

\subsection{Operational Flows}

Management is based on two types of movements performed by operators:

\begin{itemize}
  \item \textbf{PUT (Stocking):} Placing raw materials in the designated location.
  \item \textbf{PICK (Picking):} Removing components to supply production lines.
\end{itemize}

The accuracy and speed of these operations are critical: any identification error or picking delay directly impacts production line continuity. It is in this context that optimization through the Pick-to-Light system is undertaken.

\begin{figure}[H]
  \centering
  \textit{[Figure placeholder: Raw materials warehouse]}
  \caption{Raw materials warehouse of ECI Tangier.}
  \label{fig:warehouse}
\end{figure}

\section{Problem Analysis and Current Process Limitations}
\label{sec:problem-analysis}

\subsection{Problem Analysis Using the 5W2H Framework}

To rigorously frame the operational problem identified within the raw materials warehouse, the 5W2H method was applied. This structured analysis tool, whose acronym stands for Who, What, Where, When, How, Why, is widely used in industrial engineering and continuous improvement approaches to break down a problem into its fundamental dimensions, ensuring that no critical aspect is overlooked before defining a solution.

Applied to the warehouse management situation at ECI Tangier, the analysis yields the structured breakdown shown in Table~\ref{tab:5w2h}.

\begin{longtable}{p{2cm} p{3cm} p{7cm}}
  \caption{5W2H analysis framework.}\label{tab:5w2h}\\
  \toprule
  \textbf{Dimension} & \textbf{Question} & \textbf{Analysis} \\
  \midrule
  \endfirsthead
  \multicolumn{3}{c}{\tablename\ \thetable\ --- Continued} \\
  \toprule
  \textbf{Dimension} & \textbf{Question} & \textbf{Analysis} \\
  \midrule
  \endhead
  \bottomrule
  \endfoot
  Who? & Who is affected by the problem? & Warehouse operators in charge of PUT and PICK operations, as well as production supervisors who depend on the exact availability of stocks \\
  What? & What is the problem? & The absence of a real-time digitalized guidance system for raw material storage and sampling operations, leading to manual errors and traceability gaps \\
  Where? & Where does the problem occur? & Raw materials store --- storage racks, ECI Tangier production site \\
  When? & When does the problem occur? & At each PUT or PICK operation triggered by a production request, occurring daily throughout the work shifts \\
  How? & How does the problem manifest? & Operators rely on printed sheets or static Excel files to locate the materials, without any confirmation mechanism to validate the accuracy of the operation performed \\
  Why? & Why does the problem exist? & No digitalized store management system has been deployed; the existing process was designed for low volumes and has not evolved to meet current production requirements \\
\end{longtable}

This structured analysis reveals that the root cause of the problem is not operator negligence, but rather a systemic lack in the warehouse management infrastructure. The absence of real-time digital guidance creates conditions where errors are not merely possible but structurally inevitable, regardless of operator experience or attention.

\subsection{Current Process Limitations}

Before the implementation of the Pick/Put-to-Light system, raw material management was based on a largely manual process. Operators identified storage locations by looking at printed sheets or static Excel files, without any real-time digital guidance. The typical operational flow was as follows:

\begin{enumerate}
  \item Upon receipt of a production request, the operator consulted the list of references to identify the assigned rack and position.
  \item Physically moved to the location.
  \item Carried out the PUT (placing in stock) or PICK (picking) operation.
  \item Manually updated the stock records.
\end{enumerate}

Although functional for low volumes, this workflow generated increasing operational inefficiencies as production rates increased.

\begin{figure}[H]
  \centering
  \textit{[Figure placeholder: Limitation diagram and operational impact]}
  \caption{Limitation diagram and operational impact.}
  \label{fig:limitations}
\end{figure}

\section{Project Objectives and Scope}
\label{sec:objectives}

The main objective of this project is to design and implement a Pick-to-Light (PTL) system within the ECI Tangier raw materials warehouse. This system aims to digitize material handling operations, eliminate manual search errors, and provide real-time visibility into stock status.

The specific project objectives are detailed below:

\begin{itemize}
  \item \textbf{LED Guidance:} Guide operators to the exact storage position via NeoPixel LEDs controlled by ESP32 microcontrollers.
  \item \textbf{Double Confirmation:} Implement a two-step barcode validation mechanism (position scan + reference scan) to eliminate errors.
  \item \textbf{Real-Time HMI:} Develop a Human-Machine Interface (HMI) for live monitoring of warehouse status and stock levels.
  \item \textbf{Full Traceability:} Ensure complete traceability of all PUT and PICK operations with data export (Excel).
  \item \textbf{Standalone Deployment:} Deploy a standalone industrial-grade application (.exe) usable without third-party development tools.
\end{itemize}

\begin{table}[H]
  \centering
  \caption{Project objectives, indicators, and targets.}
  \label{tab:objectives}
  \begin{tabular}{p{4cm} p{5cm} p{4cm}}
    \toprule
    \textbf{Objective} & \textbf{Indicator} & \textbf{Target} \\
    \midrule
    Reduce picking errors & Error rate (\%) & Less than 0.1\% \\
    Improve picking speed & Picks per hour & +30\% improvement \\
    Real-time guidance & LED response time & Less than 500 ms \\
    Intuitive operation & Training time & Less than 30 minutes \\
    Traceability & Data logged per pick & 100\% of picks \\
    \bottomrule
  \end{tabular}
\end{table}

\section{Project Management Specifications}
\label{sec:specifications}

These specifications define the formal framework of the project, specifying the expected objectives, the constraints to be respected, the scope of intervention as well as the deliverables and success criteria.

\begin{longtable}{p{3.5cm} p{9cm}}
  \caption{Project specifications.}\label{tab:project-specs}\\
  \toprule
  \textbf{Attribute} & \textbf{Details} \\
  \midrule
  \endfirsthead
  \multicolumn{2}{c}{\tablename\ \thetable\ --- Continued} \\
  \toprule
  \textbf{Attribute} & \textbf{Details} \\
  \midrule
  \endhead
  \bottomrule
  \endfoot
  \textbf{Project Title} & Design and Implementation of a Smart Storage Management System based on Pick/Put-To-Light Technology --- ECI Tangier Raw Materials Store \\
  \textbf{Project Owner} & Mohamed D AGMOUMI ESSAFI --- Process Engineer intern, Process Engineering Department \\
  \textbf{Company Supervisor} & Mr. Saad Bouamar --- Process Engineer, ECI Tangier \\
  \textbf{Period} & February 2026 --- July 5, 2026 (5 months) \\
  \textbf{Main Objective} & Digitize the management of ECI Tangier's raw materials warehouse by deploying a low-cost Pick-to-Light system based on IoT technologies \\
  \textbf{Specific Obj. 1} & Implement a light guide using NeoPixel LEDs controlled via an ESP32 microcontroller \\
  \textbf{Specific Obj. 2} & Develop a double validation mechanism by barcode scanning (position + reference) \\
  \textbf{Specific Obj. 3} & Design a real-time HMI in Python/CustomTkinter for store supervision \\
  \textbf{Specific Obj. 4} & Ensure the complete traceability of movements with Excel export \\
  \textbf{Specific Obj. 5} & Deploy a standalone application (.exe) without technical dependencies \\
  \textbf{Specific Obj. 6} & Minimize total cost by using low-cost components and open source software \\
  \textbf{Tech Constraints} & Guided only by LEDs (no sound sensors or screens per location) \\
  \textbf{Comm Constraints} & Wireless communication via WiFi + MQTT protocol (no dedicated cabling) \\
  \textbf{Financial Constraint} & Minimal budget, low-cost components and open-source tools only \\
  \textbf{Time Constraint} & Delivery of the functional system before July 5, 2026 \\
  \textbf{Org. Constraint} & Immediate handling by operators without prior technical training \\
  \textbf{Included Scope} & 7 racks, 252 active positions, PUT and PICK operations, HMI + ESP32 + MQTT + Excel \\
  \textbf{Excluded Scope} & Other areas of the factory, multi-site management \\
  \textbf{Deliverable 1} & PTL System Hardware (ESP32 + LEDs) operational on the 7 warehouse racks \\
  \textbf{Deliverable 2} & Python HMI standalone application with real-time monitoring, KPIs and Excel file export \\
  \textbf{Deliverable 3} & Internship report complete document structured by chapters \\
  \textbf{Success Criterion 1} & The LED at the correct location lights up in less than 2 seconds after entering \\
  \textbf{Success Criterion 2} & No operation confirmed without double scan (position + reference) \\
  \textbf{Success Criterion 3} & All operations are recorded and exportable in real time \\
  \textbf{Success Criterion 4} & Functional application in .exe without installing a technical environment \\
\end{longtable}

\section{Stakeholders}
\label{sec:stakeholders}

\begin{table}[H]
  \centering
  \caption{Project stakeholders.}
  \label{tab:stakeholders}
  \begin{tabular}{p{3cm} p{5cm} p{5cm}}
    \toprule
    \textbf{Stakeholder} & \textbf{Role} & \textbf{Interest} \\
    \midrule
    Warehouse Operators & End users & Easy-to-use guidance system, reduced walking effort \\
    Warehouse Supervisor & Operational manager & Real-time tracking, error reduction reports \\
    Maintenance Team & Technical support & Reliable hardware, ease of repair \\
    IT Department & Infrastructure support & Network connectivity, PC deployment \\
    Project Sponsor & Funding & ROI, measurable performance improvement \\
    \bottomrule
  \end{tabular}
\end{table}

\section{Functional Analysis Diagrams}
\label{sec:functional-analysis}

\subsection{B\^ete \`a Cornes Diagram}
\label{sec:beteacornes}

The B\^ete \`a Cornes (Horned Beast) diagram identifies the three fundamental elements of the system: the product (SSMS), the user (Warehouse Operator), and the global goal.

\begin{figure}[H]
  \centering
  \begin{tikzpicture}[
      box/.style={rectangle, draw, rounded corners, minimum width=3.5cm, minimum height=1.5cm, align=center, font=\small, thick},
      arrow/.style={-{Latex[length=2mm]}, thick}
    ]
    \node[box] (product) at (0,0) {SSMS};
    \node[box] (user) at (6,0) {Warehouse\\Operator};
    \node[box] (goal) at (3,-2.5) {Reduce picking errors\\and improve\\retrieval speed};

    \draw[arrow] (product) -- (goal);
    \draw[arrow] (user) -- (goal);
    \draw[arrow] (product) -- (user);
  \end{tikzpicture}
  \caption{B\^ete \`a Cornes diagram for the SSMS.}
  \label{fig:beteacornes}
\end{figure}

\subsection{Pieuvre / Octopus Diagram}
\label{sec:pieuvre}

The Pieuvre (Octopus) diagram shows the main function (FP1) and the constraint functions (FC1--FC5) that link the SSMS to its environment.

\begin{figure}[H]
  \centering
  \begin{tikzpicture}[
      node distance=2.5cm,
      main/.style={ellipse, draw, minimum width=2cm, minimum height=1.2cm, align=center, fill=ssmsBlue!10, thick},
      env/.style={rectangle, draw, rounded corners, minimum width=2.5cm, minimum height=1cm, align=center, font=\footnotesize, thick},
      link/.style={-{Latex[length=2mm]}, thick}
    ]
    \node[main] (system) at (0,0) {SSMS};

    \node[env] (operator) at (0,-3) {Warehouse\\Operator};
    \node[env] (power) at (-4,-2) {Power Supply\\(230 VAC)};
    \node[env] (mqtt) at (-4,1) {MQTT\\Broker};
    \node[env] (pc) at (0,3) {Desktop\\PC};
    \node[env] (mount) at (4,1) {Rack\\Structure};
    \node[env] (safety) at (4,-2) {Safety\\Standards};

    \draw[link] (system) -- node[sloped,above,font=\tiny] {FP1} (operator);
    \draw[link] (system) -- node[sloped,above,font=\tiny] {FC1} (power);
    \draw[link] (system) -- node[sloped,above,font=\tiny] {FC2} (mqtt);
    \draw[link] (system) -- node[sloped,above,font=\tiny] {FC3} (pc);
    \draw[link] (system) -- node[sloped,above,font=\tiny] {FC4} (mount);
    \draw[link] (system) -- node[sloped,above,font=\tiny] {FC5} (safety);
  \end{tikzpicture}
  \caption{Pieuvre (Octopus) diagram for the SSMS.}
  \label{fig:pieuvre}
\end{figure}

\vspace{0.5cm}
\noindent\textbf{Function descriptions:}
\begin{itemize}
  \item \textbf{FP1 (Main Function):} Guide the warehouse operator to the correct storage bin using visual LED cues.
  \item \textbf{FC1 (Constraint):} Be powered by the mains electrical supply via the HLK-PM01 AC/DC module.
  \item \textbf{FC2 (Constraint):} Communicate with the desktop application over MQTT via the local LAN.
  \item \textbf{FC3 (Constraint):} Interface with the desktop PC for order management and data logging.
  \item \textbf{FC4 (Constraint):} Be physically mountable on the warehouse rack structure.
  \item \textbf{FC5 (Constraint):} Comply with applicable safety standards for low-voltage electrical equipment.
\end{itemize}

\section{Planning: Gantt Chart}
\label{sec:planning}

\begin{figure}[H]
  \centering
  \textit{[Figure placeholder: Gantt Chart]}
  \caption{Gantt chart for project planning (February to June 2026).}
  \label{fig:gantt}
\end{figure}

\section{Improvement Approach --- V-Cycle Methodology}
\label{sec:vcycle}

The V-Model (V-Cycle) is a widely used system development life cycle model in engineering, particularly for projects combining hardware and software. Emerging as an evolution of the waterfall model in the 1980s--1990s, especially in the aerospace, automotive, and defense sectors, it was designed to overcome the limitations of strictly linear approaches by integrating test planning from the earliest design phases rather than at the end of the project~\cite{hercog2022pick}.

In the context of this project, the V-Model structures the work into six phases:

\begin{enumerate}
  \item \textbf{Requirements Analysis ---} Gathering functional and non-functional requirements from stakeholders and establishing baseline indicators of the current system.
  \item \textbf{System Design ---} Defining the overall architecture, decomposing into subsystems, and selecting the technology stack and communication protocols.
  \item \textbf{Detailed Design ---} Fine specification of each subsystem: hardware, desktop application, MQTT layer, and ESP32 firmware.
  \item \textbf{Implementation ---} Development and programming of each component: hardware prototype, Python application, broker configuration, and MicroPython firmware.
  \item \textbf{Integration Testing ---} Verifying the joint operation of components: MQTT communication, hardware-software linkage, and end-to-end processes.
  \item \textbf{Validation \& Verification ---} Confirming that the system meets all initial requirements: error rate reduction, recovery time, scalability, and reliability.
\end{enumerate}

This approach is not presented as an isolated section of the report but as a transversal logic visible throughout the chapters: each design phase corresponds to a testing phase, ensuring that every requirement defined upstream is explicitly verified downstream.

\section{Conclusion}
\label{sec:ch1-conclusion}

This chapter has laid the contextual foundations of the project. Through the presentation of ECI Tangier and its industrial environment, we have identified the limitations of the manual warehouse management process and precisely defined the objectives of the system to be designed. The adopted methodological framework, structured around the V-Model approach, ensures a rigorous and coherent approach throughout this work.

In this approach, this chapter represents the beginning of the PLAN phase: problem analysis, needs definition, and project planning. The following chapter focuses on establishing the theoretical and technological foundations upon which the developed solution is based.
